\section{Первообразная функции на интервале. Единственность первообразной с точностью до константы. Неопределенный интеграл и его свойства.}

\subsection{Первообразная функции}

\begin{definition}
    Пусть функция \(f(x)\) определена на интервале \((a, b)\). 
    Функция \(F(x)\) называется \textbf{первообразной} для \(f(x)\) на \((a, b)\), если
    \[
    F'(x) = f(x) \quad \text{для всех } x \in (a, b).
    \]
\end{definition}

\subsection{Единственность первообразной с точностью до константы}

\begin{theorem}
    \textbf{О виде функции, производная которой равна нулю.}
    
    Если функция \(F(x)\) непрерывна на \([a, b]\) и дифференцируема на \((a, b)\), причём
    \[
    F'(x) = 0 \quad \text{для всех } x \in (a, b),
    \]
    то \(F(x) = \text{const}\) на \([a, b]\).
\end{theorem}

\begin{Proof}
    Возьмём любые \(x_1, x_2 \in [a, b]\), \(x_1 < x_2\). 
    На отрезке \([x_1, x_2]\) функция \(F\) удовлетворяет условиям теоремы Лагранжа.
    
    По теореме Лагранжа \(\exists \xi \in (x_1, x_2)\) такое, что
    \[
    F(x_2) - F(x_1) = F'(\xi)(x_2 - x_1).
    \]
    
    Но \(F'(\xi) = 0\) по условию, следовательно,
    \[
    F(x_2) - F(x_1) = 0 \quad \Rightarrow \quad F(x_2) = F(x_1).
    \]
    
    Так как \(x_1, x_2\) произвольны, функция \(F\) постоянна на \([a, b]\).
\end{Proof}

\begin{theorem}
    \textbf{Единственность первообразной.}
    
    Пусть \(F_1(x)\) и \(F_2(x)\) — первообразные для одной и той же функции \(f(x)\) на интервале \((a, b)\). Тогда
    \[
    F_1(x) - F_2(x) = \text{const} \quad \text{на } (a, b).
    \]
\end{theorem}

\begin{Proof}
    Рассмотрим \(G(x) = F_1(x) - F_2(x)\) на \((a, b)\). Тогда
    \[
    G'(x) = F_1'(x) - F_2'(x) = f(x) - f(x) = 0 \quad \text{на } (a, b).
    \]
    По теореме о функции, производная которой равна нулю, получаем \(G(x) = \text{const}\) на \((a, b)\).
\end{Proof}

\subsection{Неопределённый интеграл}

\begin{definition}
    \textbf{Неопределённым интегралом} от функции \(f(x)\) на интервале \((a, b)\) называется множество всех её первообразных.
    Обозначение:
    \[
    \int f(x) \, dx = \{ F(x) + C \mid C \in \mathbb{R} \},
    \]
    где \(F(x)\) — какая-либо первообразная для \(f(x)\).
    
    На практике пишут \(\int f(x) \, dx = F(x) + C\), понимая под \(C\) произвольную постоянную.
\end{definition}

\subsection{Свойства неопределённого интеграла}

\begin{theorem}[Свойства неопределённого интеграла]
    \begin{enumerate}
        \item \(\displaystyle \left( \int f(x) \, dx \right)' = f(x)\)
        \item \(\displaystyle d\left( \int f(x) \, dx \right) = f(x) \, dx\)
        \item \(\displaystyle \int dF(x) = F(x) + C\)
        \item \(\displaystyle \int F'(x) \, dx = F(x) + C\)
        \item \(\displaystyle \int (f(x) + g(x)) \, dx = \int f(x) \, dx + \int g(x) \, dx\) \quad (линейность)
        \item \(\displaystyle \int \alpha f(x) \, dx = \alpha \int f(x) \, dx\) \quad (линейность)
        \item \(\displaystyle \int f(x)g'(x) \, dx = f(x)g(x) - \int f'(x)g(x) \, dx\) \quad (интегрирование по частям)
    \end{enumerate}
\end{theorem}

\begin{Proof}
    \textbf{Свойство 1:} Если \(\int f(x) \, dx = F(x) + C\), то \((F(x) + C)' = F'(x) = f(x)\).
    
    \textbf{Свойство 2:} \(d\left( \int f(x) \, dx \right) = \left( \int f(x) \, dx \right)' dx = f(x) \, dx\).
    
    \textbf{Свойство 3:} По определению дифференциала, \(\int dF(x) = F(x) + C\).
    
    \textbf{Свойство 4:} \(\int F'(x) \, dx = F(x) + C\) по определению первообразной.
    
    \textbf{Свойства 5-6 (линейность, не было на лекции):} 
    \[
    \int (f(x) + g(x)) \, dx = \int f(x) \, dx + \int g(x) \, dx, \quad
    \int \alpha f(x) \, dx = \alpha \int f(x) \, dx
    \]
    следуют из линейности производной.
    
\end{Proof}

